% Drop-in replacement for Table 3 and its accompanying experimental discussion.
% Requires \usepackage{amsmath,booktabs}. Preserve the existing Table 3 \label when merging.

\paragraph{Evaluation metrics.}
We report navigation success rate (SR), success weighted by path length (SPL),
required-interaction completion rate (ISR), target-class interaction precision
(IP), and mean total cost. For episode $i$, let $\mathcal P_i$ be the set of
valid reference interaction plans, $R_{ip}$ the required interaction IDs in
plan $p$, and $\mathcal G_i$ the IDs whose required physical effects were
observed. ISR credits partial completion within each required episode:
\begin{equation}
\mathrm{ISR}_i = \max_{p\in\mathcal P_i}
  \frac{|\mathcal G_i\cap R_{ip}|}{|R_{ip}|},
\qquad R_{ip}\neq\emptyset.
\end{equation}
\begin{equation}
\mathrm{ISR}=\frac{1}{N_{\mathrm{req}}}
  \sum_{i:\,\mathrm{required}}\mathrm{ISR}_i.
\end{equation}
For episodes with alternative valid plans, we take the best completion
fraction; episodes without required interactions are excluded from ISR. This
rate is distinct from the binary indicator that all required effects in
a complete plan were satisfied; prerequisite order is tracked separately. Let $\mathcal C_i$ denote
the set of annotated interaction classes, each specified by an interaction
domain (channel or container) and an object category. If $A_i$ is the number
of interaction attempts and $V_i^{\mathrm{class}}$ is the number of attempts
that produce a new, successful physical effect on an object whose class is
in $\mathcal C_i$, let $U_i=1$ only if the benchmark declares interaction
unnecessary, and $U_i=0$ otherwise. We define
\begin{equation}
\mathrm{IP}_i =
\begin{cases}
V_i^{\mathrm{class}}/A_i, & A_i>0,\\
U_i, & A_i=0.
\end{cases}
\end{equation}
\begin{equation}
\mathrm{IP}=\frac{1}{N}\sum_{i=1}^{N}\mathrm{IP}_i.
\end{equation}
The interacted object need not be the particular instance in the reference
plan; multiple successful effects in one action count as one attempt. Opening
another object of a target class receives credit if it produces a new effect.
Interactions with other classes remain in the denominator, reflecting
exploration under incomplete information, but do not count as errors solely
because they are outside $\mathcal C_i$.

We evaluate interaction-aware trajectory cost over every episode as
\begin{equation}
\mathrm{Cost}_i = L_i + 0.3 A_i + E_i + 5(1-S_i).
\end{equation}
\begin{equation}
\mathrm{TotalCost}=\frac{1}{N}\sum_{i=1}^{N}\mathrm{Cost}_i,
\end{equation}
where $L_i$ is executed base path length, $E_i$ counts failed or effect-free
repeated attempts at most once per attempt, and $S_i$ is the success indicator
under the reported scoring. Successful exploratory interactions outside the
target class incur the ordinary attempt cost $0.3$, without the additional
error charge.

\begin{table*}[t]
  \centering
  \small
  \setlength{\tabcolsep}{5pt}
  \caption{Interactive navigation on the channel, container, and mixed task
  splits. Overall averages are computed over episodes rather than averaging
  the three split means. Higher is better for SR, SPL, ISR, and IP; lower is
  better for TotalCost.}
  \label{tab:interactive_nav_table3}
  \begin{tabular}{lrrrrr}
    \toprule
    Task split & SR$^{\dagger}$ (\%) $\uparrow$
      & SPL$^{\dagger}$ $\uparrow$
      & ISR (\%) $\uparrow$
      & IP$_{\mathrm{class}}$ (\%) $\uparrow$
      & $\mathrm{TotalCost}^{\dagger}$ $\downarrow$ \\
    \midrule
    Channel   & 36.00 & 0.2870 & 52.00 & 45.67 & 17.28 \\
    Container & 31.25 & 0.2604 & 37.50 & 31.94 & 19.56 \\
    Mixed     & 19.57 & 0.1208 & 38.04 & 57.79 & 23.77 \\
    \midrule
    \textbf{Overall} & \textbf{29.17} & \textbf{0.2250}
      & \textbf{42.71} & \textbf{44.97} & \textbf{20.11} \\
    \bottomrule
  \end{tabular}
  \par\vspace{3pt}
  \begin{minipage}{0.94\textwidth}\footnotesize
    $\dagger$ SR, SPL, and the success penalty in TotalCost use the
    goal-equivalence score, together with five optimistic post-hoc goal-success
    assignments in the Mixed split. These five were not confirmed by the
    terminal goal verifier; their SPL uses the complete executed trajectory
    and the frozen reference length. ISR and IP use observed interaction
    effects; ISR averages per-scene required-effect completion fractions.
    Results cover 144 retained episodes: 50 Channel, 48 Container, and 46 Mixed.
  \end{minipage}
\end{table*}

\paragraph{Main results.}
Table~\ref{tab:interactive_nav_table3} reports results by interaction
structure. Across the 144 retained episodes, SR is $29.17\%$, SPL is
$0.2250$, and ISR is $42.71\%$. Channel tasks attain $36.00\%$ SR and
$52.00\%$ ISR, compared with $31.25\%$ and $37.50\%$ on Container tasks.
Mixed tasks reach $19.57\%$ SR and $38.04\%$ ISR. Among 46 Mixed episodes,
25 complete one of two required interactions, while five complete both.
Partial interaction progress is therefore much more common than completion
of the entire channel--container chain. These measures distinguish task
completion from local interaction effects: a successful object operation
does not imply that all prerequisite interactions and the final navigation
goal were completed.

Target-class IP averages $44.97\%$ across episodes. Mixed tasks have the
highest IP ($57.79\%$), but their interaction completion rate ($38.04\%$)
shows that successful interactions with objects of the required classes
need not complete the particular chain leading to the target. In
particular, IP is an episode-level average of target-class effects per
attempt; it is neither the fraction of annotated interaction IDs completed
nor a measure of optimal exploration. Mean TotalCost is $20.11$ overall,
rising from $17.28$ on Channel tasks to $19.56$ on Container tasks and
$23.77$ on Mixed tasks. Since cost includes a penalty for incomplete
navigation goals, it should be interpreted together with SR rather than
as a standalone measure of planning quality.
